\documentclass[aspectratio=169,10pt]{beamer}
\makeatletter\def\input@path{{../../estilo/}{../estilo/}{estilo/}}\makeatother
\usepackage{devmobile}
\modulo[Flutter]{14}{Testes e Qualidade}{Testes de unidade, de widget, \emph{golden} e de integração, \emph{mocks}, acessibilidade, análise estática, cobertura e integração contínua}

\begin{document}
\capa

\begin{frame}{Objetivos do módulo}
  \begin{conceito}[Ao final deste módulo você será capaz de]
    \begin{itemize}
      \item Explicar a pirâmide de testes e o que testar em cada nível.
      \item Testar view models, repositórios e \emph{streams} com \emph{fakes} e \cod{mocktail}.
      \item Escrever testes de widget, de acessibilidade e \emph{golden}.
      \item Automatizar fluxos completos com \cod{integration\_test} e Patrol.
      \item Configurar análise estática, formatação e cobertura.
      \item Rodar tudo a cada \emph{push} com GitHub Actions.
    \end{itemize}
  \end{conceito}
\end{frame}

\sumario

% =============================================================================
%  Fundamentos compartilhados --- Módulo 14: teste de software
%  Incluído pelas trilhas Kotlin e Flutter com % =============================================================================
%  Fundamentos compartilhados --- Módulo 14: teste de software
%  Incluído pelas trilhas Kotlin e Flutter com % =============================================================================
%  Fundamentos compartilhados --- Módulo 14: teste de software
%  Incluído pelas trilhas Kotlin e Flutter com \input{../fundamentos/f14-testes}
% =============================================================================
\section{Fundamentos de teste de software}

\begin{frame}{Erro, defeito e falha}
  \begin{definicao}[Três palavras que não são sinônimos]
    \textbf{Erro} (\emph{error}): engano cometido por uma pessoa.
    \textbf{Defeito} (\emph{defect}, \emph{bug}): imperfeição no código ou em outro artefato, consequência do erro.
    \textbf{Falha} (\emph{failure}): comportamento incorreto observado quando o defeito é executado.
  \end{definicao}
  \centering
  \begin{tikzpicture}[node distance=17mm]
    \node[caixa,minimum width=3.0cm] (e) {\textbf{Erro}\\ \tiny o dev usa \texttt{>} em vez de \texttt{>=}};
    \node[caixapetroleo,minimum width=3.0cm,right=of e] (d) {\textbf{Defeito}\\ \tiny a regra de alerta fica errada};
    \node[caixaambar,minimum width=3.2cm,right=of d] (f) {\textbf{Falha}\\ \tiny a 30,0\,°C nenhum alerta aparece};
    \draw[seta] (e) -- node[rotulo,above]{\tiny introduz} (d);
    \draw[seta] (d) -- node[rotulo,above]{\tiny executado} (f);
  \end{tikzpicture}
  \vspace{2pt}
  \begin{columns}[T]
    \begin{column}{0.48\textwidth}
      \begin{conceito}[Consequência]
        Um defeito pode ficar anos sem causar falha, se a linha nunca for executada nas condições certas. Testar é \textbf{provocar} essas condições.
      \end{conceito}
    \end{column}
    \begin{column}{0.48\textwidth}
      \begin{historia}
        Em 1947, a equipe do computador Harvard Mark II, da qual Grace Hopper fazia parte, achou uma mariposa presa num relé e a colou no diário de bordo. A palavra \emph{bug} já existia, mas o episódio a popularizou.
      \end{historia}
    \end{column}
  \end{columns}
\end{frame}

\begin{frame}{Para que serve testar}
  \begin{columns}[T]
    \begin{column}{0.5\textwidth}
      \begin{conceito}[Objetivos dos testes]
        \begin{itemize}
          \item Encontrar defeitos antes do usuário.
          \item Dar confiança para entregar.
          \item Impedir que defeitos corrigidos voltem (\textbf{regressão}).
          \item Documentar, em código executável, como o sistema deve se comportar.
          \item Permitir \textbf{refatorar} sem medo.
        \end{itemize}
      \end{conceito}
      \begin{atencao}[Princípio fundamental]
        Testes mostram que defeitos \textbf{existem}; nunca provam que \textbf{não existem}. A ideia é de Edsger Dijkstra.
      \end{atencao}
    \end{column}
    \begin{column}{0.48\textwidth}
      \begin{pordentro}[Teste exaustivo é impossível]
        Uma função com dois inteiros de 32 bits tem $2^{64}$ combinações de entrada, cerca de $1{,}8\times10^{19}$. A um bilhão de testes por segundo, seriam quase 600 anos. Por isso escolhemos casos \textbf{representativos} com técnicas de projeto de testes.
      \end{pordentro}
      \begin{analogia}
        O \emph{crash test} de um carro não bate o veículo em todas as velocidades e ângulos: escolhe poucos cenários que representam os acidentes reais.
      \end{analogia}
    \end{column}
  \end{columns}
\end{frame}

\begin{frame}{O custo de achar tarde}
  \begin{columns}[c]
    \begin{column}{0.52\textwidth}
      \centering
      \begin{tikzpicture}[x=0.95cm,y=0.5cm]
        \draw[cinzamedio] (0,0) -- (5.3,0);
        \foreach \x/\h/\t in {0.4/0.6/Requisitos,1.4/1.2/Projeto,2.4/2.2/Código,3.4/3.8/Teste,4.4/6.4/Produção}{
          \fill[azulpetroleo] (\x,0) rectangle ++(0.6,\h);
          \node[font=\tiny,rotate=35,anchor=north east] at (\x+0.45,-0.1) {\t};}
        \fill[ambar] (4.4,0) rectangle ++(0.6,6.4);
        \node[rotulo,font=\tiny,align=center] at (2.2,6.5) {custo relativo de corrigir\\ (ilustrativo)};
      \end{tikzpicture}
    \end{column}
    \begin{column}{0.46\textwidth}
      \begin{conceito}[Por que custa mais]
        Um defeito achado em produção exige investigar relatos, reproduzir, corrigir, testar de novo, gerar versão, passar pela revisão da loja e esperar os usuários atualizarem. No app, ainda há dados já gravados errado.
      \end{conceito}
      \begin{dica}[\emph{Shift left}]
        Antecipar a qualidade: revisar requisitos, escrever testes junto com o código e rodá-los a cada \emph{commit}.
      \end{dica}
      \begin{historia}
        Barry Boehm mediu esse efeito nos anos 1970 e 1980; os números variam entre projetos, mas a tendência se repete.
      \end{historia}
    \end{column}
  \end{columns}
\end{frame}

\begin{frame}{Verificação, validação e níveis de teste}
  \begin{columns}[c]
    \begin{column}{0.56\textwidth}
      \centering
      \begin{tikzpicture}[every node/.style={font=\scriptsize,minimum width=2.25cm,minimum height=6mm}]
        \node[caixa] (l1) at (1.15,0) {Requisitos};
        \node[caixa] (l2) at (1.55,-0.75) {Especificação};
        \node[caixa] (l3) at (1.95,-1.5) {Arquitetura};
        \node[caixa] (l4) at (2.35,-2.25) {Projeto detalhado};
        \node[caixaescura] (c) at (3.8,-3.0) {Código};
        \node[caixaambar] (r4) at (5.25,-2.25) {Teste de unidade};
        \node[caixaambar] (r3) at (5.65,-1.5) {Integração};
        \node[caixaambar] (r2) at (6.05,-0.75) {Teste de sistema};
        \node[caixaambar] (r1) at (6.45,0) {Aceitação};
        \foreach \a/\b in {l1/l2,l2/l3,l3/l4,l4/c,c/r4,r4/r3,r3/r2,r2/r1} \draw[seta] (\a) -- (\b);
        \foreach \a/\b in {l1/r1,l2/r2,l3/r3,l4/r4} \draw[pontilhado,<->] (\a) -- (\b);
      \end{tikzpicture}\\[2pt]
      {\tiny Modelo em V: cada nível de teste confere um nível de especificação.}
    \end{column}
    \begin{column}{0.42\textwidth}
      \begin{definicao}[Verificação]
        ``Estamos construindo o produto \textbf{corretamente}?'' Confere o software contra a especificação.
      \end{definicao}
      \begin{definicao}[Validação]
        ``Estamos construindo o produto \textbf{certo}?'' Confere se o software resolve o problema real do usuário.
      \end{definicao}
      \begin{dica}
        O teste de aceitação é feito com o cliente: no LabMonitor, os técnicos do laboratório usando o app.
      \end{dica}
    \end{column}
  \end{columns}
\end{frame}

\begin{frame}{Caixa-preta: partições e valores-limite}
  \begin{columns}[T]
    \begin{column}{0.5\textwidth}
      \begin{definicao}[Teste caixa-preta]
        Projetado a partir da \textbf{especificação}, sem olhar o código: entradas e saídas esperadas.
      \end{definicao}
      \begin{conceito}[Partição de equivalência]
        Divide as entradas em classes que o sistema deve tratar do mesmo jeito; basta um representante de cada classe.
      \end{conceito}
      \begin{conceito}[Análise de valor-limite]
        Defeitos se concentram nas bordas (\texttt{<} no lugar de \texttt{<=}). Teste o limite e seus vizinhos imediatos.
      \end{conceito}
    \end{column}
    \begin{column}{0.48\textwidth}
      \begin{exemplo}[Regra: a reserva dura de 1 a 4 horas inteiras]
        \centering
        \begin{tikzpicture}[x=0.62cm]
          \draw[cinzamedio,-{Stealth}] (-0.5,0) -- (6.5,0);
          \fill[vermelhoerro!25] (-0.5,-0.15) rectangle (0.5,0.15);
          \fill[verdeok!35] (0.5,-0.15) rectangle (4.5,0.15);
          \fill[vermelhoerro!25] (4.5,-0.15) rectangle (6.5,0.15);
          \foreach \v in {0,1,2,3,4,5,6} \node[font=\tiny,below] at (\v,-0.15) {\v};
          \node[font=\tiny,above] at (0,0.15) {inválida};
          \node[font=\tiny,above] at (2.5,0.15) {válida};
          \node[font=\tiny,above] at (5.5,0.15) {inválida};
        \end{tikzpicture}\par
        \scriptsize
        \begin{tabular}{@{}l l@{}}
          \toprule
          \cabecalho{Casos} & \cabecalho{Por quê} \\
          \midrule
          2 h & representante da classe válida \\
          0 h e 1 h & limite inferior e vizinho \\
          4 h e 5 h & limite superior e vizinho \\
          $-1$ h, ``abc'' & entradas absurdas \\
          \bottomrule
        \end{tabular}
      \end{exemplo}
    \end{column}
  \end{columns}
\end{frame}

\begin{frame}[fragile]{Caixa-branca: critérios de cobertura}
  \begin{columns}[T]
    \begin{column}{0.46\textwidth}
      \begin{definicao}[Teste caixa-branca]
        Projetado a partir da \textbf{estrutura interna} do código: comandos, decisões e caminhos.
      \end{definicao}
\begin{codigo}[title=Pseudocódigo]{bash}
se temp >= limite E alertaAtivo:
    notificar()
registrar(temp)
\end{codigo}
      \begin{atencao}
        Cobertura mede o que foi \textbf{executado}, não o que foi \textbf{conferido}: um teste sem \cod{expect} também ``cobre''.
      \end{atencao}
    \end{column}
    \begin{column}{0.52\textwidth}
      \scriptsize
      \begin{tabularx}{\linewidth}{@{}p{2.1cm} L@{}}
        \toprule
        \cabecalho{Critério} & \cabecalho{Exige, no exemplo} \\
        \midrule
        \textbf{Comandos} & cada linha rodar ao menos uma vez: 1 teste (30\,°C, ativo) basta \\
        \textbf{Decisões} & cada \cod{se} (\emph{branch}) ser verdadeiro e falso: 2 testes \\
        \textbf{Condições} & cada condição simples verdadeira e falsa: \cod{temp} e \cod{alertaAtivo} variando \\
        \textbf{Caminhos} & todas as rotas possíveis; cresce exponencialmente com laços e \cod{se}s \\
        \bottomrule
      \end{tabularx}
      \begin{conceito}[Hierarquia]
        Caminhos $\Rightarrow$ decisões $\Rightarrow$ comandos: quem satisfaz um critério mais forte satisfaz os mais fracos.
      \end{conceito}
    \end{column}
  \end{columns}
\end{frame}

\begin{frame}{Dublês de teste}
  \begin{columns}[T]
    \begin{column}{0.62\textwidth}
      \scriptsize
      \begin{tabularx}{\linewidth}{@{}l L L@{}}
        \toprule
        \cabecalho{Dublê} & \cabecalho{O que é} & \cabecalho{Exemplo no LabMonitor} \\
        \midrule
        \textbf{Dummy} & objeto só para preencher um parâmetro; nunca é usado & um \emph{logger} vazio exigido pelo construtor \\
        \textbf{Stub} & devolve respostas prontas & repositório que sempre devolve 3 bancadas \\
        \textbf{Spy} & stub que também registra as chamadas & conta quantas notificações foram enviadas \\
        \textbf{Mock} & objeto com \textbf{expectativas} de interação, verificadas no teste & ``\cod{reservar(7)} deve ser chamado uma vez'' \\
        \textbf{Fake} & implementação funcional simplificada & banco em memória, API falsa com dados fixos \\
        \bottomrule
      \end{tabularx}
    \end{column}
    \begin{column}{0.36\textwidth}
      \begin{definicao}[Dublê de teste]
        Objeto que substitui uma dependência real (rede, banco, sensor, relógio) durante o teste.
      \end{definicao}
      \begin{analogia}
        Como o dublê de cinema: entra nas cenas perigosas ou caras no lugar do ator.
      \end{analogia}
      \begin{historia}
        A classificação é de Gerard Meszaros, no livro \emph{xUnit Test Patterns} (2007).
      \end{historia}
    \end{column}
  \end{columns}
\end{frame}

\begin{frame}{Propriedades de um bom teste}
  \begin{columns}[T]
    \begin{column}{0.5\textwidth}
      \scriptsize
      \begin{tabularx}{\linewidth}{@{}l L@{}}
        \toprule
        \cabecalho{FIRST} & \cabecalho{Significado} \\
        \midrule
        \textbf{F}ast & rápido: roda em milissegundos, então roda sempre \\
        \textbf{I}ndependent & não depende da ordem nem de outros testes \\
        \textbf{R}epeatable & mesmo resultado em qualquer máquina e horário \\
        \textbf{S}elf-validating & passa ou falha sozinho, sem alguém olhar a saída \\
        \textbf{T}imely & escrito junto com o código (ou antes) \\
        \bottomrule
      \end{tabularx}
      \begin{conceito}[Estrutura]
        \textbf{AAA}: \emph{Arrange}, \emph{Act}, \emph{Assert}.\\
        \textbf{BDD}: \emph{Given}, \emph{When}, \emph{Then} (dado, quando, então).
      \end{conceito}
    \end{column}
    \begin{column}{0.48\textwidth}
      \begin{erro}[Teste instável (\emph{flaky})]
        Passa e falha sem o código mudar. Causas típicas: depender do relógio real, da rede, da ordem de execução, de esperas fixas (\emph{sleep}) ou de estado compartilhado entre testes.
      \end{erro}
      \begin{dica}[Teste também é código]
        Merece nomes claros, nada de duplicação e revisão. Um bom nome de teste descreve uma \textbf{regra}: ``reserva sobreposta é recusada''.
      \end{dica}
    \end{column}
  \end{columns}
\end{frame}

\begin{frame}[fragile]{TDD, BDD e regressão}
  \begin{columns}[T]
    \begin{column}{0.4\textwidth}
      \centering
      \begin{tikzpicture}
        \node[circle,fill=vermelhoerro,text=white,font=\scriptsize\bfseries,minimum size=1.35cm,align=center] (v) at (90:1.35) {Vermelho\\ \tiny teste falha};
        \node[circle,fill=verdeok,text=white,font=\scriptsize\bfseries,minimum size=1.35cm,align=center] (g) at (-30:1.35) {Verde\\ \tiny passa};
        \node[circle,fill=azulpetroleo,text=white,font=\scriptsize\bfseries,minimum size=1.35cm,align=center] (r) at (210:1.35) {Refatorar\\ \tiny melhora};
        \draw[setaambar] (v) to[bend left=25] (g);
        \draw[setaambar] (g) to[bend left=25] (r);
        \draw[setaambar] (r) to[bend left=25] (v);
      \end{tikzpicture}
      \begin{definicao}[TDD]
        \emph{Test-Driven Development} (Kent Beck): o teste é escrito \textbf{antes} do código e guia o projeto.
      \end{definicao}
    \end{column}
    \begin{column}{0.58\textwidth}
\begin{codigo}[title=BDD: cenário em linguagem natural (Gherkin)]{bash}
Funcionalidade: Reserva de bancada
  Cenário: horário já ocupado
    Dado que a bancada B7 está reservada das 14h às 16h
    Quando eu peço a B7 das 15h às 17h
    Então vejo "Horário indisponível"
\end{codigo}
      \begin{definicao}[BDD]
        \emph{Behavior-Driven Development}: descreve o comportamento em exemplos que clientes, testadores e devs entendem; os cenários viram testes.
      \end{definicao}
      \begin{definicao}[Teste de regressão]
        Rodar de novo os testes existentes após cada mudança, para garantir que nada que funcionava quebrou. Automatizar é o que torna isso viável.
      \end{definicao}
    \end{column}
  \end{columns}
\end{frame}

\begin{frame}{Desafios do teste em apps móveis}
  \begin{columns}[T]
    \begin{column}{0.5\textwidth}
      \scriptsize
      \begin{tabularx}{\linewidth}{@{}p{2.0cm} L@{}}
        \toprule
        \cabecalho{Desafio} & \cabecalho{O que pode dar errado} \\
        \midrule
        Fragmentação & milhares de modelos Android, versões de sistema, tamanhos e densidades de tela \\
        Rede & sem sinal, 3G lento, troca Wi-Fi/móvel no meio de uma requisição \\
        Interrupções & ligação, notificação, app enviado ao fundo e encerrado pelo sistema \\
        Hardware & câmera, GPS, Bluetooth e sensores diferentes ou ausentes \\
        Permissões & usuário nega, concede ``só desta vez'' ou revoga depois \\
        Recursos & bateria fraca, pouca memória, armazenamento cheio \\
        Pessoas & idioma, fonte grande, leitor de tela, modo escuro \\
        \bottomrule
      \end{tabularx}
    \end{column}
    \begin{column}{0.48\textwidth}
      \begin{conceito}[Estratégia]
        \begin{itemize}
          \item Automatizar a lógica em testes de unidade e widget (rápidos, sem aparelho).
          \item Definir uma \textbf{matriz de dispositivos} representativa (versões antigas e novas, tela pequena e grande).
          \item Usar emuladores no dia a dia e aparelhos reais ou \emph{device farms} antes de publicar.
          \item Simular rede ruim e interrupções.
        \end{itemize}
      \end{conceito}
      \begin{dica}
        As estatísticas do Android Studio e do console da loja mostram quais versões seus usuários realmente usam.
      \end{dica}
    \end{column}
  \end{columns}
\end{frame}

\begin{frame}{Qualidade é mais do que testar}
  \begin{columns}[T]
    \begin{column}{0.5\textwidth}
      \begin{definicao}[Modelo de qualidade ISO/IEC 25010]
        Organiza a qualidade de um produto de software em características como adequação funcional, eficiência de desempenho, compatibilidade, usabilidade, confiabilidade, segurança, manutenibilidade e portabilidade.
      \end{definicao}
      \begin{dica}
        A revisão de 2023 renomeou algumas características (usabilidade virou capacidade de interação) e incluiu \emph{safety}.
      \end{dica}
    \end{column}
    \begin{column}{0.48\textwidth}
      \begin{conceito}[Técnicas que se somam]
        \begin{itemize}
          \item \textbf{Estáticas} (sem executar): revisão de código, análise estática, \emph{lint}, tipagem forte.
          \item \textbf{Dinâmicas} (executando): testes automatizados e exploratórios.
          \item \textbf{Em produção}: monitoramento de falhas, métricas e \emph{feedback} dos usuários.
        \end{itemize}
      \end{conceito}
      \begin{destaque}
        Qualidade se \textbf{constrói} durante o desenvolvimento; testes confirmam que ela está lá.
      \end{destaque}
    \end{column}
  \end{columns}
\end{frame}

% =============================================================================
\section{Fundamentos de teste de software}

\begin{frame}{Erro, defeito e falha}
  \begin{definicao}[Três palavras que não são sinônimos]
    \textbf{Erro} (\emph{error}): engano cometido por uma pessoa.
    \textbf{Defeito} (\emph{defect}, \emph{bug}): imperfeição no código ou em outro artefato, consequência do erro.
    \textbf{Falha} (\emph{failure}): comportamento incorreto observado quando o defeito é executado.
  \end{definicao}
  \centering
  \begin{tikzpicture}[node distance=17mm]
    \node[caixa,minimum width=3.0cm] (e) {\textbf{Erro}\\ \tiny o dev usa \texttt{>} em vez de \texttt{>=}};
    \node[caixapetroleo,minimum width=3.0cm,right=of e] (d) {\textbf{Defeito}\\ \tiny a regra de alerta fica errada};
    \node[caixaambar,minimum width=3.2cm,right=of d] (f) {\textbf{Falha}\\ \tiny a 30,0\,°C nenhum alerta aparece};
    \draw[seta] (e) -- node[rotulo,above]{\tiny introduz} (d);
    \draw[seta] (d) -- node[rotulo,above]{\tiny executado} (f);
  \end{tikzpicture}
  \vspace{2pt}
  \begin{columns}[T]
    \begin{column}{0.48\textwidth}
      \begin{conceito}[Consequência]
        Um defeito pode ficar anos sem causar falha, se a linha nunca for executada nas condições certas. Testar é \textbf{provocar} essas condições.
      \end{conceito}
    \end{column}
    \begin{column}{0.48\textwidth}
      \begin{historia}
        Em 1947, a equipe do computador Harvard Mark II, da qual Grace Hopper fazia parte, achou uma mariposa presa num relé e a colou no diário de bordo. A palavra \emph{bug} já existia, mas o episódio a popularizou.
      \end{historia}
    \end{column}
  \end{columns}
\end{frame}

\begin{frame}{Para que serve testar}
  \begin{columns}[T]
    \begin{column}{0.5\textwidth}
      \begin{conceito}[Objetivos dos testes]
        \begin{itemize}
          \item Encontrar defeitos antes do usuário.
          \item Dar confiança para entregar.
          \item Impedir que defeitos corrigidos voltem (\textbf{regressão}).
          \item Documentar, em código executável, como o sistema deve se comportar.
          \item Permitir \textbf{refatorar} sem medo.
        \end{itemize}
      \end{conceito}
      \begin{atencao}[Princípio fundamental]
        Testes mostram que defeitos \textbf{existem}; nunca provam que \textbf{não existem}. A ideia é de Edsger Dijkstra.
      \end{atencao}
    \end{column}
    \begin{column}{0.48\textwidth}
      \begin{pordentro}[Teste exaustivo é impossível]
        Uma função com dois inteiros de 32 bits tem $2^{64}$ combinações de entrada, cerca de $1{,}8\times10^{19}$. A um bilhão de testes por segundo, seriam quase 600 anos. Por isso escolhemos casos \textbf{representativos} com técnicas de projeto de testes.
      \end{pordentro}
      \begin{analogia}
        O \emph{crash test} de um carro não bate o veículo em todas as velocidades e ângulos: escolhe poucos cenários que representam os acidentes reais.
      \end{analogia}
    \end{column}
  \end{columns}
\end{frame}

\begin{frame}{O custo de achar tarde}
  \begin{columns}[c]
    \begin{column}{0.52\textwidth}
      \centering
      \begin{tikzpicture}[x=0.95cm,y=0.5cm]
        \draw[cinzamedio] (0,0) -- (5.3,0);
        \foreach \x/\h/\t in {0.4/0.6/Requisitos,1.4/1.2/Projeto,2.4/2.2/Código,3.4/3.8/Teste,4.4/6.4/Produção}{
          \fill[azulpetroleo] (\x,0) rectangle ++(0.6,\h);
          \node[font=\tiny,rotate=35,anchor=north east] at (\x+0.45,-0.1) {\t};}
        \fill[ambar] (4.4,0) rectangle ++(0.6,6.4);
        \node[rotulo,font=\tiny,align=center] at (2.2,6.5) {custo relativo de corrigir\\ (ilustrativo)};
      \end{tikzpicture}
    \end{column}
    \begin{column}{0.46\textwidth}
      \begin{conceito}[Por que custa mais]
        Um defeito achado em produção exige investigar relatos, reproduzir, corrigir, testar de novo, gerar versão, passar pela revisão da loja e esperar os usuários atualizarem. No app, ainda há dados já gravados errado.
      \end{conceito}
      \begin{dica}[\emph{Shift left}]
        Antecipar a qualidade: revisar requisitos, escrever testes junto com o código e rodá-los a cada \emph{commit}.
      \end{dica}
      \begin{historia}
        Barry Boehm mediu esse efeito nos anos 1970 e 1980; os números variam entre projetos, mas a tendência se repete.
      \end{historia}
    \end{column}
  \end{columns}
\end{frame}

\begin{frame}{Verificação, validação e níveis de teste}
  \begin{columns}[c]
    \begin{column}{0.56\textwidth}
      \centering
      \begin{tikzpicture}[every node/.style={font=\scriptsize,minimum width=2.25cm,minimum height=6mm}]
        \node[caixa] (l1) at (1.15,0) {Requisitos};
        \node[caixa] (l2) at (1.55,-0.75) {Especificação};
        \node[caixa] (l3) at (1.95,-1.5) {Arquitetura};
        \node[caixa] (l4) at (2.35,-2.25) {Projeto detalhado};
        \node[caixaescura] (c) at (3.8,-3.0) {Código};
        \node[caixaambar] (r4) at (5.25,-2.25) {Teste de unidade};
        \node[caixaambar] (r3) at (5.65,-1.5) {Integração};
        \node[caixaambar] (r2) at (6.05,-0.75) {Teste de sistema};
        \node[caixaambar] (r1) at (6.45,0) {Aceitação};
        \foreach \a/\b in {l1/l2,l2/l3,l3/l4,l4/c,c/r4,r4/r3,r3/r2,r2/r1} \draw[seta] (\a) -- (\b);
        \foreach \a/\b in {l1/r1,l2/r2,l3/r3,l4/r4} \draw[pontilhado,<->] (\a) -- (\b);
      \end{tikzpicture}\\[2pt]
      {\tiny Modelo em V: cada nível de teste confere um nível de especificação.}
    \end{column}
    \begin{column}{0.42\textwidth}
      \begin{definicao}[Verificação]
        ``Estamos construindo o produto \textbf{corretamente}?'' Confere o software contra a especificação.
      \end{definicao}
      \begin{definicao}[Validação]
        ``Estamos construindo o produto \textbf{certo}?'' Confere se o software resolve o problema real do usuário.
      \end{definicao}
      \begin{dica}
        O teste de aceitação é feito com o cliente: no LabMonitor, os técnicos do laboratório usando o app.
      \end{dica}
    \end{column}
  \end{columns}
\end{frame}

\begin{frame}{Caixa-preta: partições e valores-limite}
  \begin{columns}[T]
    \begin{column}{0.5\textwidth}
      \begin{definicao}[Teste caixa-preta]
        Projetado a partir da \textbf{especificação}, sem olhar o código: entradas e saídas esperadas.
      \end{definicao}
      \begin{conceito}[Partição de equivalência]
        Divide as entradas em classes que o sistema deve tratar do mesmo jeito; basta um representante de cada classe.
      \end{conceito}
      \begin{conceito}[Análise de valor-limite]
        Defeitos se concentram nas bordas (\texttt{<} no lugar de \texttt{<=}). Teste o limite e seus vizinhos imediatos.
      \end{conceito}
    \end{column}
    \begin{column}{0.48\textwidth}
      \begin{exemplo}[Regra: a reserva dura de 1 a 4 horas inteiras]
        \centering
        \begin{tikzpicture}[x=0.62cm]
          \draw[cinzamedio,-{Stealth}] (-0.5,0) -- (6.5,0);
          \fill[vermelhoerro!25] (-0.5,-0.15) rectangle (0.5,0.15);
          \fill[verdeok!35] (0.5,-0.15) rectangle (4.5,0.15);
          \fill[vermelhoerro!25] (4.5,-0.15) rectangle (6.5,0.15);
          \foreach \v in {0,1,2,3,4,5,6} \node[font=\tiny,below] at (\v,-0.15) {\v};
          \node[font=\tiny,above] at (0,0.15) {inválida};
          \node[font=\tiny,above] at (2.5,0.15) {válida};
          \node[font=\tiny,above] at (5.5,0.15) {inválida};
        \end{tikzpicture}\par
        \scriptsize
        \begin{tabular}{@{}l l@{}}
          \toprule
          \cabecalho{Casos} & \cabecalho{Por quê} \\
          \midrule
          2 h & representante da classe válida \\
          0 h e 1 h & limite inferior e vizinho \\
          4 h e 5 h & limite superior e vizinho \\
          $-1$ h, ``abc'' & entradas absurdas \\
          \bottomrule
        \end{tabular}
      \end{exemplo}
    \end{column}
  \end{columns}
\end{frame}

\begin{frame}[fragile]{Caixa-branca: critérios de cobertura}
  \begin{columns}[T]
    \begin{column}{0.46\textwidth}
      \begin{definicao}[Teste caixa-branca]
        Projetado a partir da \textbf{estrutura interna} do código: comandos, decisões e caminhos.
      \end{definicao}
\begin{codigo}[title=Pseudocódigo]{bash}
se temp >= limite E alertaAtivo:
    notificar()
registrar(temp)
\end{codigo}
      \begin{atencao}
        Cobertura mede o que foi \textbf{executado}, não o que foi \textbf{conferido}: um teste sem \cod{expect} também ``cobre''.
      \end{atencao}
    \end{column}
    \begin{column}{0.52\textwidth}
      \scriptsize
      \begin{tabularx}{\linewidth}{@{}p{2.1cm} L@{}}
        \toprule
        \cabecalho{Critério} & \cabecalho{Exige, no exemplo} \\
        \midrule
        \textbf{Comandos} & cada linha rodar ao menos uma vez: 1 teste (30\,°C, ativo) basta \\
        \textbf{Decisões} & cada \cod{se} (\emph{branch}) ser verdadeiro e falso: 2 testes \\
        \textbf{Condições} & cada condição simples verdadeira e falsa: \cod{temp} e \cod{alertaAtivo} variando \\
        \textbf{Caminhos} & todas as rotas possíveis; cresce exponencialmente com laços e \cod{se}s \\
        \bottomrule
      \end{tabularx}
      \begin{conceito}[Hierarquia]
        Caminhos $\Rightarrow$ decisões $\Rightarrow$ comandos: quem satisfaz um critério mais forte satisfaz os mais fracos.
      \end{conceito}
    \end{column}
  \end{columns}
\end{frame}

\begin{frame}{Dublês de teste}
  \begin{columns}[T]
    \begin{column}{0.62\textwidth}
      \scriptsize
      \begin{tabularx}{\linewidth}{@{}l L L@{}}
        \toprule
        \cabecalho{Dublê} & \cabecalho{O que é} & \cabecalho{Exemplo no LabMonitor} \\
        \midrule
        \textbf{Dummy} & objeto só para preencher um parâmetro; nunca é usado & um \emph{logger} vazio exigido pelo construtor \\
        \textbf{Stub} & devolve respostas prontas & repositório que sempre devolve 3 bancadas \\
        \textbf{Spy} & stub que também registra as chamadas & conta quantas notificações foram enviadas \\
        \textbf{Mock} & objeto com \textbf{expectativas} de interação, verificadas no teste & ``\cod{reservar(7)} deve ser chamado uma vez'' \\
        \textbf{Fake} & implementação funcional simplificada & banco em memória, API falsa com dados fixos \\
        \bottomrule
      \end{tabularx}
    \end{column}
    \begin{column}{0.36\textwidth}
      \begin{definicao}[Dublê de teste]
        Objeto que substitui uma dependência real (rede, banco, sensor, relógio) durante o teste.
      \end{definicao}
      \begin{analogia}
        Como o dublê de cinema: entra nas cenas perigosas ou caras no lugar do ator.
      \end{analogia}
      \begin{historia}
        A classificação é de Gerard Meszaros, no livro \emph{xUnit Test Patterns} (2007).
      \end{historia}
    \end{column}
  \end{columns}
\end{frame}

\begin{frame}{Propriedades de um bom teste}
  \begin{columns}[T]
    \begin{column}{0.5\textwidth}
      \scriptsize
      \begin{tabularx}{\linewidth}{@{}l L@{}}
        \toprule
        \cabecalho{FIRST} & \cabecalho{Significado} \\
        \midrule
        \textbf{F}ast & rápido: roda em milissegundos, então roda sempre \\
        \textbf{I}ndependent & não depende da ordem nem de outros testes \\
        \textbf{R}epeatable & mesmo resultado em qualquer máquina e horário \\
        \textbf{S}elf-validating & passa ou falha sozinho, sem alguém olhar a saída \\
        \textbf{T}imely & escrito junto com o código (ou antes) \\
        \bottomrule
      \end{tabularx}
      \begin{conceito}[Estrutura]
        \textbf{AAA}: \emph{Arrange}, \emph{Act}, \emph{Assert}.\\
        \textbf{BDD}: \emph{Given}, \emph{When}, \emph{Then} (dado, quando, então).
      \end{conceito}
    \end{column}
    \begin{column}{0.48\textwidth}
      \begin{erro}[Teste instável (\emph{flaky})]
        Passa e falha sem o código mudar. Causas típicas: depender do relógio real, da rede, da ordem de execução, de esperas fixas (\emph{sleep}) ou de estado compartilhado entre testes.
      \end{erro}
      \begin{dica}[Teste também é código]
        Merece nomes claros, nada de duplicação e revisão. Um bom nome de teste descreve uma \textbf{regra}: ``reserva sobreposta é recusada''.
      \end{dica}
    \end{column}
  \end{columns}
\end{frame}

\begin{frame}[fragile]{TDD, BDD e regressão}
  \begin{columns}[T]
    \begin{column}{0.4\textwidth}
      \centering
      \begin{tikzpicture}
        \node[circle,fill=vermelhoerro,text=white,font=\scriptsize\bfseries,minimum size=1.35cm,align=center] (v) at (90:1.35) {Vermelho\\ \tiny teste falha};
        \node[circle,fill=verdeok,text=white,font=\scriptsize\bfseries,minimum size=1.35cm,align=center] (g) at (-30:1.35) {Verde\\ \tiny passa};
        \node[circle,fill=azulpetroleo,text=white,font=\scriptsize\bfseries,minimum size=1.35cm,align=center] (r) at (210:1.35) {Refatorar\\ \tiny melhora};
        \draw[setaambar] (v) to[bend left=25] (g);
        \draw[setaambar] (g) to[bend left=25] (r);
        \draw[setaambar] (r) to[bend left=25] (v);
      \end{tikzpicture}
      \begin{definicao}[TDD]
        \emph{Test-Driven Development} (Kent Beck): o teste é escrito \textbf{antes} do código e guia o projeto.
      \end{definicao}
    \end{column}
    \begin{column}{0.58\textwidth}
\begin{codigo}[title=BDD: cenário em linguagem natural (Gherkin)]{bash}
Funcionalidade: Reserva de bancada
  Cenário: horário já ocupado
    Dado que a bancada B7 está reservada das 14h às 16h
    Quando eu peço a B7 das 15h às 17h
    Então vejo "Horário indisponível"
\end{codigo}
      \begin{definicao}[BDD]
        \emph{Behavior-Driven Development}: descreve o comportamento em exemplos que clientes, testadores e devs entendem; os cenários viram testes.
      \end{definicao}
      \begin{definicao}[Teste de regressão]
        Rodar de novo os testes existentes após cada mudança, para garantir que nada que funcionava quebrou. Automatizar é o que torna isso viável.
      \end{definicao}
    \end{column}
  \end{columns}
\end{frame}

\begin{frame}{Desafios do teste em apps móveis}
  \begin{columns}[T]
    \begin{column}{0.5\textwidth}
      \scriptsize
      \begin{tabularx}{\linewidth}{@{}p{2.0cm} L@{}}
        \toprule
        \cabecalho{Desafio} & \cabecalho{O que pode dar errado} \\
        \midrule
        Fragmentação & milhares de modelos Android, versões de sistema, tamanhos e densidades de tela \\
        Rede & sem sinal, 3G lento, troca Wi-Fi/móvel no meio de uma requisição \\
        Interrupções & ligação, notificação, app enviado ao fundo e encerrado pelo sistema \\
        Hardware & câmera, GPS, Bluetooth e sensores diferentes ou ausentes \\
        Permissões & usuário nega, concede ``só desta vez'' ou revoga depois \\
        Recursos & bateria fraca, pouca memória, armazenamento cheio \\
        Pessoas & idioma, fonte grande, leitor de tela, modo escuro \\
        \bottomrule
      \end{tabularx}
    \end{column}
    \begin{column}{0.48\textwidth}
      \begin{conceito}[Estratégia]
        \begin{itemize}
          \item Automatizar a lógica em testes de unidade e widget (rápidos, sem aparelho).
          \item Definir uma \textbf{matriz de dispositivos} representativa (versões antigas e novas, tela pequena e grande).
          \item Usar emuladores no dia a dia e aparelhos reais ou \emph{device farms} antes de publicar.
          \item Simular rede ruim e interrupções.
        \end{itemize}
      \end{conceito}
      \begin{dica}
        As estatísticas do Android Studio e do console da loja mostram quais versões seus usuários realmente usam.
      \end{dica}
    \end{column}
  \end{columns}
\end{frame}

\begin{frame}{Qualidade é mais do que testar}
  \begin{columns}[T]
    \begin{column}{0.5\textwidth}
      \begin{definicao}[Modelo de qualidade ISO/IEC 25010]
        Organiza a qualidade de um produto de software em características como adequação funcional, eficiência de desempenho, compatibilidade, usabilidade, confiabilidade, segurança, manutenibilidade e portabilidade.
      \end{definicao}
      \begin{dica}
        A revisão de 2023 renomeou algumas características (usabilidade virou capacidade de interação) e incluiu \emph{safety}.
      \end{dica}
    \end{column}
    \begin{column}{0.48\textwidth}
      \begin{conceito}[Técnicas que se somam]
        \begin{itemize}
          \item \textbf{Estáticas} (sem executar): revisão de código, análise estática, \emph{lint}, tipagem forte.
          \item \textbf{Dinâmicas} (executando): testes automatizados e exploratórios.
          \item \textbf{Em produção}: monitoramento de falhas, métricas e \emph{feedback} dos usuários.
        \end{itemize}
      \end{conceito}
      \begin{destaque}
        Qualidade se \textbf{constrói} durante o desenvolvimento; testes confirmam que ela está lá.
      \end{destaque}
    \end{column}
  \end{columns}
\end{frame}

% =============================================================================
\section{Fundamentos de teste de software}

\begin{frame}{Erro, defeito e falha}
  \begin{definicao}[Três palavras que não são sinônimos]
    \textbf{Erro} (\emph{error}): engano cometido por uma pessoa.
    \textbf{Defeito} (\emph{defect}, \emph{bug}): imperfeição no código ou em outro artefato, consequência do erro.
    \textbf{Falha} (\emph{failure}): comportamento incorreto observado quando o defeito é executado.
  \end{definicao}
  \centering
  \begin{tikzpicture}[node distance=17mm]
    \node[caixa,minimum width=3.0cm] (e) {\textbf{Erro}\\ \tiny o dev usa \texttt{>} em vez de \texttt{>=}};
    \node[caixapetroleo,minimum width=3.0cm,right=of e] (d) {\textbf{Defeito}\\ \tiny a regra de alerta fica errada};
    \node[caixaambar,minimum width=3.2cm,right=of d] (f) {\textbf{Falha}\\ \tiny a 30,0\,°C nenhum alerta aparece};
    \draw[seta] (e) -- node[rotulo,above]{\tiny introduz} (d);
    \draw[seta] (d) -- node[rotulo,above]{\tiny executado} (f);
  \end{tikzpicture}
  \vspace{2pt}
  \begin{columns}[T]
    \begin{column}{0.48\textwidth}
      \begin{conceito}[Consequência]
        Um defeito pode ficar anos sem causar falha, se a linha nunca for executada nas condições certas. Testar é \textbf{provocar} essas condições.
      \end{conceito}
    \end{column}
    \begin{column}{0.48\textwidth}
      \begin{historia}
        Em 1947, a equipe do computador Harvard Mark II, da qual Grace Hopper fazia parte, achou uma mariposa presa num relé e a colou no diário de bordo. A palavra \emph{bug} já existia, mas o episódio a popularizou.
      \end{historia}
    \end{column}
  \end{columns}
\end{frame}

\begin{frame}{Para que serve testar}
  \begin{columns}[T]
    \begin{column}{0.5\textwidth}
      \begin{conceito}[Objetivos dos testes]
        \begin{itemize}
          \item Encontrar defeitos antes do usuário.
          \item Dar confiança para entregar.
          \item Impedir que defeitos corrigidos voltem (\textbf{regressão}).
          \item Documentar, em código executável, como o sistema deve se comportar.
          \item Permitir \textbf{refatorar} sem medo.
        \end{itemize}
      \end{conceito}
      \begin{atencao}[Princípio fundamental]
        Testes mostram que defeitos \textbf{existem}; nunca provam que \textbf{não existem}. A ideia é de Edsger Dijkstra.
      \end{atencao}
    \end{column}
    \begin{column}{0.48\textwidth}
      \begin{pordentro}[Teste exaustivo é impossível]
        Uma função com dois inteiros de 32 bits tem $2^{64}$ combinações de entrada, cerca de $1{,}8\times10^{19}$. A um bilhão de testes por segundo, seriam quase 600 anos. Por isso escolhemos casos \textbf{representativos} com técnicas de projeto de testes.
      \end{pordentro}
      \begin{analogia}
        O \emph{crash test} de um carro não bate o veículo em todas as velocidades e ângulos: escolhe poucos cenários que representam os acidentes reais.
      \end{analogia}
    \end{column}
  \end{columns}
\end{frame}

\begin{frame}{O custo de achar tarde}
  \begin{columns}[c]
    \begin{column}{0.52\textwidth}
      \centering
      \begin{tikzpicture}[x=0.95cm,y=0.5cm]
        \draw[cinzamedio] (0,0) -- (5.3,0);
        \foreach \x/\h/\t in {0.4/0.6/Requisitos,1.4/1.2/Projeto,2.4/2.2/Código,3.4/3.8/Teste,4.4/6.4/Produção}{
          \fill[azulpetroleo] (\x,0) rectangle ++(0.6,\h);
          \node[font=\tiny,rotate=35,anchor=north east] at (\x+0.45,-0.1) {\t};}
        \fill[ambar] (4.4,0) rectangle ++(0.6,6.4);
        \node[rotulo,font=\tiny,align=center] at (2.2,6.5) {custo relativo de corrigir\\ (ilustrativo)};
      \end{tikzpicture}
    \end{column}
    \begin{column}{0.46\textwidth}
      \begin{conceito}[Por que custa mais]
        Um defeito achado em produção exige investigar relatos, reproduzir, corrigir, testar de novo, gerar versão, passar pela revisão da loja e esperar os usuários atualizarem. No app, ainda há dados já gravados errado.
      \end{conceito}
      \begin{dica}[\emph{Shift left}]
        Antecipar a qualidade: revisar requisitos, escrever testes junto com o código e rodá-los a cada \emph{commit}.
      \end{dica}
      \begin{historia}
        Barry Boehm mediu esse efeito nos anos 1970 e 1980; os números variam entre projetos, mas a tendência se repete.
      \end{historia}
    \end{column}
  \end{columns}
\end{frame}

\begin{frame}{Verificação, validação e níveis de teste}
  \begin{columns}[c]
    \begin{column}{0.56\textwidth}
      \centering
      \begin{tikzpicture}[every node/.style={font=\scriptsize,minimum width=2.25cm,minimum height=6mm}]
        \node[caixa] (l1) at (1.15,0) {Requisitos};
        \node[caixa] (l2) at (1.55,-0.75) {Especificação};
        \node[caixa] (l3) at (1.95,-1.5) {Arquitetura};
        \node[caixa] (l4) at (2.35,-2.25) {Projeto detalhado};
        \node[caixaescura] (c) at (3.8,-3.0) {Código};
        \node[caixaambar] (r4) at (5.25,-2.25) {Teste de unidade};
        \node[caixaambar] (r3) at (5.65,-1.5) {Integração};
        \node[caixaambar] (r2) at (6.05,-0.75) {Teste de sistema};
        \node[caixaambar] (r1) at (6.45,0) {Aceitação};
        \foreach \a/\b in {l1/l2,l2/l3,l3/l4,l4/c,c/r4,r4/r3,r3/r2,r2/r1} \draw[seta] (\a) -- (\b);
        \foreach \a/\b in {l1/r1,l2/r2,l3/r3,l4/r4} \draw[pontilhado,<->] (\a) -- (\b);
      \end{tikzpicture}\\[2pt]
      {\tiny Modelo em V: cada nível de teste confere um nível de especificação.}
    \end{column}
    \begin{column}{0.42\textwidth}
      \begin{definicao}[Verificação]
        ``Estamos construindo o produto \textbf{corretamente}?'' Confere o software contra a especificação.
      \end{definicao}
      \begin{definicao}[Validação]
        ``Estamos construindo o produto \textbf{certo}?'' Confere se o software resolve o problema real do usuário.
      \end{definicao}
      \begin{dica}
        O teste de aceitação é feito com o cliente: no LabMonitor, os técnicos do laboratório usando o app.
      \end{dica}
    \end{column}
  \end{columns}
\end{frame}

\begin{frame}{Caixa-preta: partições e valores-limite}
  \begin{columns}[T]
    \begin{column}{0.5\textwidth}
      \begin{definicao}[Teste caixa-preta]
        Projetado a partir da \textbf{especificação}, sem olhar o código: entradas e saídas esperadas.
      \end{definicao}
      \begin{conceito}[Partição de equivalência]
        Divide as entradas em classes que o sistema deve tratar do mesmo jeito; basta um representante de cada classe.
      \end{conceito}
      \begin{conceito}[Análise de valor-limite]
        Defeitos se concentram nas bordas (\texttt{<} no lugar de \texttt{<=}). Teste o limite e seus vizinhos imediatos.
      \end{conceito}
    \end{column}
    \begin{column}{0.48\textwidth}
      \begin{exemplo}[Regra: a reserva dura de 1 a 4 horas inteiras]
        \centering
        \begin{tikzpicture}[x=0.62cm]
          \draw[cinzamedio,-{Stealth}] (-0.5,0) -- (6.5,0);
          \fill[vermelhoerro!25] (-0.5,-0.15) rectangle (0.5,0.15);
          \fill[verdeok!35] (0.5,-0.15) rectangle (4.5,0.15);
          \fill[vermelhoerro!25] (4.5,-0.15) rectangle (6.5,0.15);
          \foreach \v in {0,1,2,3,4,5,6} \node[font=\tiny,below] at (\v,-0.15) {\v};
          \node[font=\tiny,above] at (0,0.15) {inválida};
          \node[font=\tiny,above] at (2.5,0.15) {válida};
          \node[font=\tiny,above] at (5.5,0.15) {inválida};
        \end{tikzpicture}\par
        \scriptsize
        \begin{tabular}{@{}l l@{}}
          \toprule
          \cabecalho{Casos} & \cabecalho{Por quê} \\
          \midrule
          2 h & representante da classe válida \\
          0 h e 1 h & limite inferior e vizinho \\
          4 h e 5 h & limite superior e vizinho \\
          $-1$ h, ``abc'' & entradas absurdas \\
          \bottomrule
        \end{tabular}
      \end{exemplo}
    \end{column}
  \end{columns}
\end{frame}

\begin{frame}[fragile]{Caixa-branca: critérios de cobertura}
  \begin{columns}[T]
    \begin{column}{0.46\textwidth}
      \begin{definicao}[Teste caixa-branca]
        Projetado a partir da \textbf{estrutura interna} do código: comandos, decisões e caminhos.
      \end{definicao}
\begin{codigo}[title=Pseudocódigo]{bash}
se temp >= limite E alertaAtivo:
    notificar()
registrar(temp)
\end{codigo}
      \begin{atencao}
        Cobertura mede o que foi \textbf{executado}, não o que foi \textbf{conferido}: um teste sem \cod{expect} também ``cobre''.
      \end{atencao}
    \end{column}
    \begin{column}{0.52\textwidth}
      \scriptsize
      \begin{tabularx}{\linewidth}{@{}p{2.1cm} L@{}}
        \toprule
        \cabecalho{Critério} & \cabecalho{Exige, no exemplo} \\
        \midrule
        \textbf{Comandos} & cada linha rodar ao menos uma vez: 1 teste (30\,°C, ativo) basta \\
        \textbf{Decisões} & cada \cod{se} (\emph{branch}) ser verdadeiro e falso: 2 testes \\
        \textbf{Condições} & cada condição simples verdadeira e falsa: \cod{temp} e \cod{alertaAtivo} variando \\
        \textbf{Caminhos} & todas as rotas possíveis; cresce exponencialmente com laços e \cod{se}s \\
        \bottomrule
      \end{tabularx}
      \begin{conceito}[Hierarquia]
        Caminhos $\Rightarrow$ decisões $\Rightarrow$ comandos: quem satisfaz um critério mais forte satisfaz os mais fracos.
      \end{conceito}
    \end{column}
  \end{columns}
\end{frame}

\begin{frame}{Dublês de teste}
  \begin{columns}[T]
    \begin{column}{0.62\textwidth}
      \scriptsize
      \begin{tabularx}{\linewidth}{@{}l L L@{}}
        \toprule
        \cabecalho{Dublê} & \cabecalho{O que é} & \cabecalho{Exemplo no LabMonitor} \\
        \midrule
        \textbf{Dummy} & objeto só para preencher um parâmetro; nunca é usado & um \emph{logger} vazio exigido pelo construtor \\
        \textbf{Stub} & devolve respostas prontas & repositório que sempre devolve 3 bancadas \\
        \textbf{Spy} & stub que também registra as chamadas & conta quantas notificações foram enviadas \\
        \textbf{Mock} & objeto com \textbf{expectativas} de interação, verificadas no teste & ``\cod{reservar(7)} deve ser chamado uma vez'' \\
        \textbf{Fake} & implementação funcional simplificada & banco em memória, API falsa com dados fixos \\
        \bottomrule
      \end{tabularx}
    \end{column}
    \begin{column}{0.36\textwidth}
      \begin{definicao}[Dublê de teste]
        Objeto que substitui uma dependência real (rede, banco, sensor, relógio) durante o teste.
      \end{definicao}
      \begin{analogia}
        Como o dublê de cinema: entra nas cenas perigosas ou caras no lugar do ator.
      \end{analogia}
      \begin{historia}
        A classificação é de Gerard Meszaros, no livro \emph{xUnit Test Patterns} (2007).
      \end{historia}
    \end{column}
  \end{columns}
\end{frame}

\begin{frame}{Propriedades de um bom teste}
  \begin{columns}[T]
    \begin{column}{0.5\textwidth}
      \scriptsize
      \begin{tabularx}{\linewidth}{@{}l L@{}}
        \toprule
        \cabecalho{FIRST} & \cabecalho{Significado} \\
        \midrule
        \textbf{F}ast & rápido: roda em milissegundos, então roda sempre \\
        \textbf{I}ndependent & não depende da ordem nem de outros testes \\
        \textbf{R}epeatable & mesmo resultado em qualquer máquina e horário \\
        \textbf{S}elf-validating & passa ou falha sozinho, sem alguém olhar a saída \\
        \textbf{T}imely & escrito junto com o código (ou antes) \\
        \bottomrule
      \end{tabularx}
      \begin{conceito}[Estrutura]
        \textbf{AAA}: \emph{Arrange}, \emph{Act}, \emph{Assert}.\\
        \textbf{BDD}: \emph{Given}, \emph{When}, \emph{Then} (dado, quando, então).
      \end{conceito}
    \end{column}
    \begin{column}{0.48\textwidth}
      \begin{erro}[Teste instável (\emph{flaky})]
        Passa e falha sem o código mudar. Causas típicas: depender do relógio real, da rede, da ordem de execução, de esperas fixas (\emph{sleep}) ou de estado compartilhado entre testes.
      \end{erro}
      \begin{dica}[Teste também é código]
        Merece nomes claros, nada de duplicação e revisão. Um bom nome de teste descreve uma \textbf{regra}: ``reserva sobreposta é recusada''.
      \end{dica}
    \end{column}
  \end{columns}
\end{frame}

\begin{frame}[fragile]{TDD, BDD e regressão}
  \begin{columns}[T]
    \begin{column}{0.4\textwidth}
      \centering
      \begin{tikzpicture}
        \node[circle,fill=vermelhoerro,text=white,font=\scriptsize\bfseries,minimum size=1.35cm,align=center] (v) at (90:1.35) {Vermelho\\ \tiny teste falha};
        \node[circle,fill=verdeok,text=white,font=\scriptsize\bfseries,minimum size=1.35cm,align=center] (g) at (-30:1.35) {Verde\\ \tiny passa};
        \node[circle,fill=azulpetroleo,text=white,font=\scriptsize\bfseries,minimum size=1.35cm,align=center] (r) at (210:1.35) {Refatorar\\ \tiny melhora};
        \draw[setaambar] (v) to[bend left=25] (g);
        \draw[setaambar] (g) to[bend left=25] (r);
        \draw[setaambar] (r) to[bend left=25] (v);
      \end{tikzpicture}
      \begin{definicao}[TDD]
        \emph{Test-Driven Development} (Kent Beck): o teste é escrito \textbf{antes} do código e guia o projeto.
      \end{definicao}
    \end{column}
    \begin{column}{0.58\textwidth}
\begin{codigo}[title=BDD: cenário em linguagem natural (Gherkin)]{bash}
Funcionalidade: Reserva de bancada
  Cenário: horário já ocupado
    Dado que a bancada B7 está reservada das 14h às 16h
    Quando eu peço a B7 das 15h às 17h
    Então vejo "Horário indisponível"
\end{codigo}
      \begin{definicao}[BDD]
        \emph{Behavior-Driven Development}: descreve o comportamento em exemplos que clientes, testadores e devs entendem; os cenários viram testes.
      \end{definicao}
      \begin{definicao}[Teste de regressão]
        Rodar de novo os testes existentes após cada mudança, para garantir que nada que funcionava quebrou. Automatizar é o que torna isso viável.
      \end{definicao}
    \end{column}
  \end{columns}
\end{frame}

\begin{frame}{Desafios do teste em apps móveis}
  \begin{columns}[T]
    \begin{column}{0.5\textwidth}
      \scriptsize
      \begin{tabularx}{\linewidth}{@{}p{2.0cm} L@{}}
        \toprule
        \cabecalho{Desafio} & \cabecalho{O que pode dar errado} \\
        \midrule
        Fragmentação & milhares de modelos Android, versões de sistema, tamanhos e densidades de tela \\
        Rede & sem sinal, 3G lento, troca Wi-Fi/móvel no meio de uma requisição \\
        Interrupções & ligação, notificação, app enviado ao fundo e encerrado pelo sistema \\
        Hardware & câmera, GPS, Bluetooth e sensores diferentes ou ausentes \\
        Permissões & usuário nega, concede ``só desta vez'' ou revoga depois \\
        Recursos & bateria fraca, pouca memória, armazenamento cheio \\
        Pessoas & idioma, fonte grande, leitor de tela, modo escuro \\
        \bottomrule
      \end{tabularx}
    \end{column}
    \begin{column}{0.48\textwidth}
      \begin{conceito}[Estratégia]
        \begin{itemize}
          \item Automatizar a lógica em testes de unidade e widget (rápidos, sem aparelho).
          \item Definir uma \textbf{matriz de dispositivos} representativa (versões antigas e novas, tela pequena e grande).
          \item Usar emuladores no dia a dia e aparelhos reais ou \emph{device farms} antes de publicar.
          \item Simular rede ruim e interrupções.
        \end{itemize}
      \end{conceito}
      \begin{dica}
        As estatísticas do Android Studio e do console da loja mostram quais versões seus usuários realmente usam.
      \end{dica}
    \end{column}
  \end{columns}
\end{frame}

\begin{frame}{Qualidade é mais do que testar}
  \begin{columns}[T]
    \begin{column}{0.5\textwidth}
      \begin{definicao}[Modelo de qualidade ISO/IEC 25010]
        Organiza a qualidade de um produto de software em características como adequação funcional, eficiência de desempenho, compatibilidade, usabilidade, confiabilidade, segurança, manutenibilidade e portabilidade.
      \end{definicao}
      \begin{dica}
        A revisão de 2023 renomeou algumas características (usabilidade virou capacidade de interação) e incluiu \emph{safety}.
      \end{dica}
    \end{column}
    \begin{column}{0.48\textwidth}
      \begin{conceito}[Técnicas que se somam]
        \begin{itemize}
          \item \textbf{Estáticas} (sem executar): revisão de código, análise estática, \emph{lint}, tipagem forte.
          \item \textbf{Dinâmicas} (executando): testes automatizados e exploratórios.
          \item \textbf{Em produção}: monitoramento de falhas, métricas e \emph{feedback} dos usuários.
        \end{itemize}
      \end{conceito}
      \begin{destaque}
        Qualidade se \textbf{constrói} durante o desenvolvimento; testes confirmam que ela está lá.
      \end{destaque}
    \end{column}
  \end{columns}
\end{frame}


% =============================================================================
\section{Testes no Flutter: visão geral}

\begin{frame}{A pirâmide de testes no Flutter}
  \begin{columns}[c]
    \begin{column}{0.44\textwidth}
      \centering
      \begin{tikzpicture}[scale=0.9,transform shape]
        \fill[ambar!80] (-1.0,2.6) -- (1.0,2.6) -- (0.45,3.5) -- (-0.45,3.5) -- cycle;
        \fill[azulprofundo!75] (-1.7,1.4) -- (1.7,1.4) -- (1.0,2.5) -- (-1.0,2.5) -- cycle;
        \fill[azulpetroleo] (-2.5,0) -- (2.5,0) -- (1.75,1.3) -- (-1.75,1.3) -- cycle;
        \node[font=\scriptsize\bfseries] at (0,3.05) {Integração};
        \node[font=\scriptsize\bfseries,text=white] at (0,1.95) {Widget};
        \node[font=\scriptsize\bfseries,text=white] at (0,0.65) {Unidade};
        \draw[setaambar] (2.9,0.1) -- node[rotulo,right,align=left]{\tiny mais lento\\ \tiny mais caro\\ \tiny mais realista} (2.9,3.4);
      \end{tikzpicture}
    \end{column}
    \begin{column}{0.54\textwidth}
      \scriptsize
      \begin{tabularx}{\linewidth}{@{}l L p{1.75cm}@{}}
        \toprule
        \cabecalho{Tipo} & \cabecalho{Testa} & \cabecalho{Roda em} \\
        \midrule
        \textbf{Unidade} & uma função, classe, view model ou repositório & VM Dart, ms \\
        \textbf{Widget} & uma tela ou componente, sem aparelho & ambiente simulado, ms \\
        \textbf{Golden} & a aparência, pixel a pixel & ambiente simulado \\
        \textbf{Integração} & o app inteiro em um fluxo real & aparelho ou emulador, s--min \\
        \bottomrule
      \end{tabularx}
      \begin{dica}[Proporção saudável]
        Muitos testes de unidade, bastantes de widget e poucos de integração, cobrindo os fluxos críticos (login, reserva).
      \end{dica}
    \end{column}
  \end{columns}
\end{frame}

% =============================================================================
\section{Testes de unidade}

\begin{frame}[fragile]{Estrutura de um teste}
  \begin{columns}[T]
    \begin{column}{0.54\textwidth}
\begin{codigo}[title=test/utils/conversao\_test.dart]{Dart}
import 'package:flutter_test/flutter_test.dart';
import 'package:labmonitor/utils/conversao.dart';

void main() {
  group('celsiusDeBytes', () {
    test('converte centésimos little-endian', () {
      expect(celsiusDeBytes([0xB7, 0x09]), 24.87);
    });

    test('aceita temperatura negativa', () {
      expect(celsiusDeBytes([0x18, 0xFC]), -10.0);
    });

    test('rejeita pacote curto', () {
      expect(() => celsiusDeBytes([0x01]),
          throwsA(isA<FormatException>()));
    });
  });
}
\end{codigo}
    \end{column}
    \begin{column}{0.44\textwidth}
\begin{terminal}
$ flutter test
00:02 +3: All tests passed!
$ flutter test test/utils/
$ flutter test --name "negativa"
\end{terminal}
      \begin{conceito}[Padrão AAA]
        \textbf{Arrange} (preparar), \textbf{Act} (agir), \textbf{Assert} (conferir). Um comportamento por teste, com nome que descreve a regra.
      \end{conceito}
      \begin{dica}[\emph{Matchers} úteis]
        \cod{equals}, \cod{isTrue}, \cod{hasLength}, \cod{contains}, \cod{closeTo(24.9, 0.1)}, \cod{throwsA}, \cod{isA<T>()}.
      \end{dica}
    \end{column}
  \end{columns}
\end{frame}

\begin{frame}[fragile]{Testando o view model com um \emph{fake}}
  \begin{columns}[T]
    \begin{column}{0.58\textwidth}
\begin{codigo}[title=test/ui/painel\_view\_model\_test.dart]{Dart}
class RepoFake implements BancadaRepository {
  RepoFake(this.resposta);
  Result<List<Bancada>> resposta;
  @override
  Future<Result<List<Bancada>>> listar() async => resposta;
}
void main() {
  final duas = Ok([bancada('B1'), bancada('B2')]);
  test('carrega ao criar e filtra pelo nome', () async {
    final vm = PainelViewModel(repo: RepoFake(duas));
    await Future<void>.delayed(Duration.zero);   // fim do load
    vm.filtrar('B2');
    expect(vm.visiveis.single.nome, 'B2');
  });
  test('expõe a falha do repositório', () async {
    final repo = RepoFake(Falha(Exception('offline')));
    final vm = PainelViewModel(repo: repo);
    await Future<void>.delayed(Duration.zero);
    expect(vm.carregar.falhou, isTrue);
  });
}
\end{codigo}
    \end{column}
    \begin{column}{0.4\textwidth}
      \begin{conceito}[\emph{Fake}]
        Implementação simples e funcional da interface, escrita por você. Controla exatamente o que o repositório devolve.
      \end{conceito}
      \begin{dica}
        A arquitetura do Módulo 08 se paga aqui: o view model recebe a interface, e o teste dispensa rede, banco e Firebase.
      \end{dica}
      \begin{atencao}
        Teste comportamento observável (estado exposto), não detalhes internos.
      \end{atencao}
    \end{column}
  \end{columns}
\end{frame}

\begin{frame}[fragile]{\emph{Mocks} com \cod{mocktail}}
  \begin{columns}[T]
    \begin{column}{0.58\textwidth}
\begin{codigo}[title=Verificando interações]{Dart}
import 'package:mocktail/mocktail.dart';

class RepoMock extends Mock implements BancadaRepository {}

void main() {
  test('tentar de novo chama o repositório outra vez',
      () async {
    final repo = RepoMock();
    when(() => repo.listar()).thenAnswer(
        (_) async => Falha(Exception('offline')));

    final vm = PainelViewModel(repo: repo);
    await Future<void>.delayed(Duration.zero);
    await vm.carregar.execute();

    verify(() => repo.listar()).called(2);
  });
}
\end{codigo}
    \end{column}
    \begin{column}{0.4\textwidth}
      \begin{conceito}[\emph{Fake} ou \emph{mock}?]
        \textbf{Fake}: quando importa o \emph{resultado}.\\
        \textbf{Mock}: quando importa a \emph{interação} (foi chamado? quantas vezes? com quais argumentos?).
      \end{conceito}
      \begin{dica}
        \cod{mocktail} não precisa de geração de código. Para argumentos com \cod{any()} de tipos próprios, chame \cod{registerFallbackValue} no \cod{setUpAll}.
      \end{dica}
    \end{column}
  \end{columns}
\end{frame}

\begin{frame}[fragile]{\emph{Streams} e tempo}
  \begin{columns}[T]
    \begin{column}{0.5\textwidth}
\begin{codigo}[title=Streams]{Dart}
test('converte os pacotes do BLE', () {
  final pacotes = Stream.fromIterable([
    [0xB7, 0x09],
    [0x00, 0x0A],
  ]);
  expect(pacotes.map(celsiusDeBytes),
      emitsInOrder([24.87, 25.6, emitsDone]));
});
\end{codigo}
      \begin{conceito}[\emph{Stream matchers}]
        \cod{emits}, \cod{emitsInOrder}, \cod{emitsThrough}, \cod{emitsError} e \cod{emitsDone} descrevem a sequência esperada.
      \end{conceito}
    \end{column}
    \begin{column}{0.48\textwidth}
\begin{codigo}[title=Tempo controlado (fake\_async)]{Dart}
test('busca espera parar de digitar', () {
  fakeAsync((tempo) {
    final repo = RepoContador();
    final vm = BuscaViewModel(repo: repo);
    vm.digitar('B');
    vm.digitar('B7');
    tempo.elapse(
        const Duration(milliseconds: 299));
    expect(repo.buscas, 0);
    tempo.elapse(
        const Duration(milliseconds: 1));
    expect(repo.buscas, 1);
  });
});
\end{codigo}
      \begin{dica}
        Nada de \cod{sleep}: o tempo avança na hora, e o teste fica rápido e determinístico.
      \end{dica}
    \end{column}
  \end{columns}
\end{frame}

% =============================================================================
\section{Testes de widget}

\begin{frame}[fragile]{Testando uma tela}
  \begin{columns}[T]
    \begin{column}{0.58\textwidth}
\begin{codigo}[title=test/ui/painel\_screen\_test.dart]{Dart}
testWidgets('mostra o erro e tenta de novo', (tester) async {
  final repo = RepoFake(Falha(Exception('offline')));
  final vm = PainelViewModel(repo: repo);

  await tester.pumpWidget(
      MaterialApp(home: PainelScreen(vm: vm)));
  await tester.pump();                  // conclui o load
  expect(find.text('Tentar de novo'), findsOneWidget);

  repo.resposta = Ok([bancada('B1')]);
  await tester.tap(find.text('Tentar de novo'));
  await tester.pump();

  expect(find.text('B1'), findsOneWidget);
  expect(find.text('Tentar de novo'), findsNothing);
});
\end{codigo}
    \end{column}
    \begin{column}{0.4\textwidth}
      \begin{conceito}[\cod{WidgetTester}]
        Monta a árvore de widgets sem aparelho. \cod{pump} desenha um quadro; \cod{pumpAndSettle} desenha até as animações terminarem.
      \end{conceito}
      \begin{atencao}
        O tempo é simulado: \cod{Future}s e \cod{Timer}s só avançam com \cod{pump}. Esquecer o \cod{pump} é a causa mais comum de teste falhando.
      \end{atencao}
    \end{column}
  \end{columns}
\end{frame}

\begin{frame}[fragile]{Encontrar, agir e conferir}
  \begin{columns}[T]
    \begin{column}{0.5\textwidth}
      \scriptsize
      \begin{tabularx}{\linewidth}{@{}l L@{}}
        \toprule
        \cabecalho{Uso} & \cabecalho{API} \\
        \midrule
        Encontrar & \cod{find.text}, \cod{find.byType}, \cod{find.byKey}, \cod{find.byIcon}, \cod{find.bySemanticsLabel} \\
        Agir & \cod{tap}, \cod{longPress}, \cod{enterText}, \cod{drag}, \cod{scrollUntilVisible} \\
        Conferir & \cod{findsOneWidget}, \cod{findsNothing}, \cod{findsNWidgets(3)}, \cod{findsAtLeast(1)} \\
        \bottomrule
      \end{tabularx}
      \begin{dica}
        Dê \cod{Key}s estáveis aos campos importantes (\cod{Key('email')}): o teste não quebra quando o texto muda.
      \end{dica}
    \end{column}
    \begin{column}{0.48\textwidth}
\begin{codigo}[title=Acessibilidade automática]{Dart}
testWidgets('painel acessível', (tester) async {
  final h = tester.ensureSemantics();
  await tester.pumpWidget(appDeTeste());
  await expectLater(tester,
      meetsGuideline(androidTapTargetGuideline));
  await expectLater(tester,
      meetsGuideline(labeledTapTargetGuideline));
  await expectLater(tester,
      meetsGuideline(textContrastGuideline));
  h.dispose();
});
\end{codigo}
      \begin{conceito}[O que é verificado]
        Alvos de toque com 48\,dp, botões com rótulo para leitores de tela e contraste de texto.
      \end{conceito}
    \end{column}
  \end{columns}
\end{frame}

% =============================================================================
\section{Testes \emph{golden}}

\begin{frame}[fragile]{Comparando imagens}
  \begin{columns}[T]
    \begin{column}{0.56\textwidth}
\begin{codigo}[title=test/ui/cartao\_bancada\_golden\_test.dart]{Dart}
testWidgets('cartão de bancada ocupada', (tester) async {
  await tester.pumpWidget(MaterialApp(
    theme: temaLab,
    home: Center(child: CartaoBancada(
        bancada: bancada('B7', livre: false))),
  ));
  await expectLater(find.byType(CartaoBancada),
      matchesGoldenFile('goldens/cartao_ocupada.png'));
});
\end{codigo}
\begin{terminal}
$ flutter test --update-goldens   # gera/atualiza
$ flutter test                    # compara
\end{terminal}
    \end{column}
    \begin{column}{0.42\textwidth}
      \begin{conceito}[\emph{Golden test}]
        Renderiza o widget, salva uma imagem de referência e, nas execuções seguintes, falha se algum pixel mudar. Pega regressões visuais que nenhum \cod{expect} pegaria.
      \end{conceito}
      \begin{atencao}
        Fontes e antialiasing variam entre sistemas: gere e compare os \emph{goldens} no mesmo ambiente (de preferência no CI).
      \end{atencao}
      \begin{dica}
        Teste os estados: claro e escuro, texto grande, carregando, erro.
      \end{dica}
    \end{column}
  \end{columns}
\end{frame}

% =============================================================================
\section{Testes de integração}

\begin{frame}[fragile]{O app inteiro no aparelho}
  \begin{columns}[T]
    \begin{column}{0.58\textwidth}
\begin{codigo}[title=integration\_test/reserva\_test.dart]{Dart}
import 'package:integration_test/integration_test.dart';
import 'package:labmonitor/main.dart' as app;

void main() {
  IntegrationTestWidgetsFlutterBinding.ensureInitialized();

  testWidgets('aluno reserva uma bancada', (tester) async {
    app.main();                 // aponta para os emuladores
    await tester.pumpAndSettle();
    await tester.enterText(find.byKey(const Key('email')),
        'aluno@teste.dev');
    await tester.enterText(find.byKey(const Key('senha')),
        'segredo123');
    await tester.tap(find.text('Entrar'));
    await tester.pumpAndSettle();
    await tester.tap(find.text('B7'));
    await tester.tap(find.text('Reservar'));
    await tester.pumpAndSettle();
    expect(find.text('Reserva confirmada'), findsOneWidget);
  });
}
\end{codigo}
    \end{column}
    \begin{column}{0.4\textwidth}
\begin{terminal}
$ flutter test integration_test
\end{terminal}
      \begin{conceito}[Mesmo jeito de escrever]
        A API é a do teste de widget, mas roda no aparelho real ou emulador, com plugins, rede e banco de verdade.
      \end{conceito}
      \begin{atencao}[Dados de teste]
        Nunca use o backend de produção. Use o Emulator Suite (Módulo 12) ou uma API de homologação.
      \end{atencao}
    \end{column}
  \end{columns}
\end{frame}

\begin{frame}[fragile]{Além do Flutter: Patrol e nuvem de aparelhos}
  \begin{columns}[T]
    \begin{column}{0.52\textwidth}
\begin{codigo}[title=Patrol: interagindo com o sistema]{Dart}
patrolTest('lê QR code após permitir a câmera',
    ($) async {
  await $.pumpWidgetAndSettle(const LabMonitorApp());
  await $('Ler QR code').tap();
  // diálogo nativo do Android/iOS
  if (await $.platform.mobile
      .isPermissionDialogVisible()) {
    await $.platform.mobile.grantPermissionWhenInUse();
  }
  await $(MobileScanner).waitUntilVisible();
});
\end{codigo}
    \end{column}
    \begin{column}{0.46\textwidth}
      \begin{conceito}[Limite do \cod{integration\_test}]
        Ele só enxerga widgets Flutter. Diálogos de permissão, notificações e configurações do sistema ficam de fora.
      \end{conceito}
      \begin{dica}[Patrol]
        Estende os testes com interação nativa e uma sintaxe mais curta de \emph{finders}.
      \end{dica}
      \begin{dica}[Firebase Test Lab]
        Roda os testes de integração em dezenas de modelos reais de Android e iOS na nuvem.
      \end{dica}
    \end{column}
  \end{columns}
\end{frame}

% =============================================================================
\section{Qualidade contínua}

\begin{frame}[fragile]{Análise estática, formatação e cobertura}
  \begin{columns}[T]
    \begin{column}{0.48\textwidth}
\begin{codigo}[title=analysis\_options.yaml]{YAML}
include: package:flutter_lints/flutter.yaml

analyzer:
  language:
    strict-casts: true
    strict-raw-types: true
  errors:
    unawaited_futures: error

linter:
  rules:
    - prefer_final_locals
    - avoid_print
    - unawaited_futures
\end{codigo}
    \end{column}
    \begin{column}{0.5\textwidth}
\begin{terminal}
$ dart format .
$ flutter analyze
No issues found!
$ flutter test --coverage
$ genhtml coverage/lcov.info -o coverage/html
\end{terminal}
      \begin{conceito}[Três guardas]
        \textbf{Formatador}: estilo único, sem discussão.\\
        \textbf{Analisador}: erros e más práticas antes de rodar.\\
        \textbf{Cobertura}: quais linhas os testes executam.
      \end{conceito}
      \begin{atencao}
        Cobertura alta não garante bons testes. Use-a para achar código \textbf{sem} teste, não como meta.
      \end{atencao}
    \end{column}
  \end{columns}
\end{frame}

\begin{frame}[fragile]{Integração contínua com GitHub Actions}
  \begin{columns}[T]
    \begin{column}{0.52\textwidth}
\begin{codigo}[title=.github/workflows/qualidade.yml]{YAML}
name: qualidade
on: [push, pull_request]

jobs:
  testes:
    runs-on: ubuntu-latest
    steps:
      - uses: actions/checkout@v5
      - uses: subosito/flutter-action@v2
        with:
          channel: stable
          cache: true
      - run: flutter pub get
      - run: dart format --set-exit-if-changed .
      - run: flutter analyze
      - run: flutter test --coverage
\end{codigo}
    \end{column}
    \begin{column}{0.46\textwidth}
      \begin{conceito}[Integração contínua]
        A cada \emph{push} e \emph{pull request}, um servidor limpo baixa o código, formata, analisa e testa. Se algo falha, o PR fica vermelho e não deve ser integrado.
      \end{conceito}
      \begin{dica}
        Proteja o \emph{branch} \cod{main}: exija o \emph{check} verde para aceitar o PR (Módulo 03).
      \end{dica}
      \begin{dica}
        O Módulo 18 estende o pipeline com \emph{build} e publicação automática.
      \end{dica}
    \end{column}
  \end{columns}
\end{frame}

\begin{frame}{Estratégia: o que testar primeiro}
  \begin{columns}[T]
    \begin{column}{0.5\textwidth}
      \begin{conceito}[Onde investir primeiro]
        Comece pelo que mais dói quando quebra e é mais barato de testar: regras de negócio e conversões de dados. Telas e fluxos completos vêm depois, cobrindo os caminhos críticos.
      \end{conceito}
      \begin{dica}[Bug corrigido = teste novo]
        Antes de corrigir um bug, escreva o teste que o reproduz. Ele nunca mais volta.
      \end{dica}
    \end{column}
    \begin{column}{0.48\textwidth}
      \scriptsize
      \begin{tabularx}{\linewidth}{@{}l L@{}}
        \toprule
        \cabecalho{Prioridade} & \cabecalho{Exemplos no LabMonitor} \\
        \midrule
        Regras de negócio & conflito de reserva, limite de alerta \\
        Conversões & bytes do BLE, JSON da API \\
        View models & carregando, erro, vazio, filtro \\
        Telas críticas & login, reserva, alertas \\
        Fluxos ponta a ponta & entrar $\rightarrow$ reservar $\rightarrow$ cancelar \\
        \bottomrule
      \end{tabularx}
      \begin{atencao}
        Não teste o framework (se o \cod{Text} mostra texto) nem pacotes de terceiros.
      \end{atencao}
    \end{column}
  \end{columns}
\end{frame}

% =============================================================================
\section{Fechamento}

\begin{frame}{Resumo}
  \begin{resumo}
    \begin{itemize}
      \item Pirâmide: muitos testes de unidade, bastantes de widget, poucos de integração.
      \item \emph{Fakes} para controlar resultados; \cod{mocktail} para verificar interações.
      \item \cod{WidgetTester} monta telas sem aparelho; \emph{goldens} protegem a aparência.
      \item \cod{integration\_test} e Patrol exercitam o app real; Test Lab multiplica aparelhos.
      \item Formatador, analisador, cobertura e CI mantêm a qualidade a cada \emph{push}.
    \end{itemize}
  \end{resumo}
\end{frame}

\begin{frame}{Glossário (1/2)}
  \begin{glossario}
    Teste de unidade & verifica uma unidade isolada (função, classe, view model), sem dependências reais \\
    Teste de widget & monta um widget ou tela em ambiente simulado e interage com ele \\
    Teste \emph{golden} & compara a renderização com uma imagem de referência salva \\
    Teste de integração & executa o app completo em aparelho ou emulador \\
    Teste de regressão & reexecução dos testes após mudanças para garantir que nada quebrou \\
    Dublê de teste & substituto de uma dependência real: \emph{dummy}, \emph{stub}, \emph{spy}, \emph{mock} ou \emph{fake} \\
    \emph{Fake} & implementação simplificada e funcional (repositório em memória) \\
    \emph{Mock} & dublê que verifica as interações esperadas (chamadas e argumentos) \\
  \end{glossario}
\end{frame}

\begin{frame}{Glossário (2/2)}
  \begin{glossario}
    \emph{Matcher} & objeto que descreve o valor esperado em um \cod{expect} (\cod{equals}, \cod{throwsA}) \\
    \emph{Finder} & localiza widgets na árvore durante o teste (\cod{find.text}, \cod{find.byKey}) \\
    \cod{WidgetTester} & controla o teste de widget: monta, desenha quadros e simula toques \\
    \cod{pump} / \cod{pumpAndSettle} & avança um quadro / avança até não haver animações pendentes \\
    Cobertura & proporção do código executada pelos testes (comandos, decisões, condições) \\
    Análise estática & inspeção do código sem executá-lo (\cod{flutter analyze}, \emph{lints}) \\
    Teste instável (\emph{flaky}) & teste que passa ou falha sem mudança no código \\
    CI & integração contínua: compilar e testar automaticamente a cada \emph{push} \\
  \end{glossario}
\end{frame}

\begin{frame}{Questões de revisão}
  \begin{questoes}
    \item Diferencie erro, defeito e falha usando um exemplo do LabMonitor.
    \item Por que testes não provam a ausência de defeitos? Que técnicas reduzem o número de casos sem perder poder de detecção?
    \item A regra é ``alerta quando a temperatura for maior ou igual a 30\,°C''. Quais valores você testaria e por quê?
    \item Mostre um exemplo em que 100\% de cobertura de comandos não implica 100\% de cobertura de decisões.
    \item Quando usar um \emph{fake} e quando usar um \emph{mock}?
    \item O que um teste de widget não consegue verificar e que exige um teste de integração?
    \item Que cuidado os testes \emph{golden} exigem no CI?
    \item Cite três causas de testes instáveis e como evitá-las.
  \end{questoes}
\end{frame}

\begin{frame}{Exercícios}
  \begin{exercicio}[Exercício 1 --- Unidade]
    Escreva testes para a regra ``uma bancada não pode ter duas reservas sobrepostas'', incluindo os limites de horário.
  \end{exercicio}
  \begin{exercicio}[Exercício 2 --- Widget e golden]
    Teste os estados carregando, vazio, erro e lista do painel e crie \emph{goldens} nos temas claro e escuro.
  \end{exercicio}
  \begin{exercicio}[Exercício 3 --- CI]
    Crie o \emph{workflow} do GitHub Actions e torne-o obrigatório para \emph{pull requests}.
  \end{exercicio}
  \begin{projeto}[LabMonitor --- Etapa 14]
    Testes dos view models e conversões, testes de widget das telas principais, um \emph{golden} por componente e um teste de integração do fluxo de reserva no CI.
  \end{projeto}
\end{frame}

\begin{frame}{Para aprofundar}
  \small
  \begin{itemize}
    \item Testes no Flutter: \url{https://docs.flutter.dev/testing/overview}
    \item Testes de integração: \url{https://docs.flutter.dev/testing/integration-tests}
    \item mocktail: \url{https://pub.dev/packages/mocktail}
    \item Patrol: \url{https://patrol.leancode.co}
    \item Regras do analisador: \url{https://dart.dev/tools/linter-rules}
    \item Firebase Test Lab: \url{https://firebase.google.com/docs/test-lab}
  \end{itemize}
\end{frame}

\begin{frame}{Respostas das questões de revisão}
  \begin{respostas}
    \item \textbf{Erro}: o dev escreve \texttt{>} em vez de \texttt{>=}. \textbf{Defeito}: a regra de alerta fica errada no código. \textbf{Falha}: a exatamente 30\,°C o alerta não aparece.
    \item O espaço de entradas é grande demais para testar tudo; um teste só mostra o comportamento nos casos executados. Partição de equivalência e valores-limite escolhem casos representativos.
    \item 29,9 (logo abaixo, sem alerta), 30,0 (no limite, com alerta) e 30,1 (acima), além de um valor típico de cada classe e entradas inválidas (sem leitura).
    \item \texttt{se cond: x()} seguido de \texttt{y()}: um teste com \texttt{cond} verdadeira executa todas as linhas, mas o lado falso da decisão nunca é exercitado.
    \item \emph{Fake} quando importa o resultado (dados devolvidos); \emph{mock} quando importa a interação (se e como uma dependência foi chamada).
    \item Plugins nativos, permissões, rede e banco reais, desempenho no aparelho e o fluxo entre telas com o app inteiro.
    \item Gerar e comparar as imagens no mesmo ambiente (sistema, fontes, versão do Flutter), de preferência no próprio CI, para evitar diferenças de renderização.
    \item Relógio real, rede real e esperas fixas; ordem de execução e estado compartilhado. Evite com \emph{fakes}, tempo simulado (\cod{fakeAsync}, \cod{pump}) e \cod{setUp} que recria o estado a cada teste.
  \end{respostas}
\end{frame}

\encerramento{Módulo 15 --- Publicação: Google Play e App Store.}

\end{document}
